\section{Introduction}
\label{sec:intro}

Vision-language-action models (VLAs) have become the dominant policy
architecture for generalist robotic manipulation, with systems such as
\textbf{$\pi_0$} \cite{pi0}, \textbf{$\pi_{0.5}$} \cite{pi05},
\textbf{OpenVLA} \cite{openvla}, \textbf{RDT-1B} \cite{rdt1b} and
\textbf{$\pi_{0.7}$} \cite{pi07} demonstrating impressive zero-shot
generalization across embodiments and tasks. However, these models are
almost universally \emph{Markovian}: each action is predicted solely
from the current observation (plus a very short observation window),
leading to systematic failure on tasks that require reasoning over
longer temporal dependencies. The RMBench benchmark
\cite{rmbench} demonstrates this starkly: mainstream VLAs achieve
near-zero success on non-Markovian tasks such as pressing a button
a variable number of times or sequencing multi-step subtasks.

A wave of \emph{memory-aware} VLAs has emerged in 2025--2026 to
address this gap:
Mem-0 \cite{rmbench} maintains a triple of anchor, sliding, and key
memories with a dedicated planner-executor hierarchy;
MemoryVLA \cite{memoryvla} introduces a perceptual-cognitive dual bank;
ReMem-VLA \cite{rememvla} proposes dual-level recurrent queries;
SAM2Act+ \cite{sam2act} adapts the SAM2 video memory for spatial
tasks; MemER \cite{memer} scales retrieval through VLM-selected
keyframes; CronusVLA \cite{cronusvla} employs multi-frame feature
chunking; and MEM \cite{mem} (the component powering
$\pi_{0.7}$) combines a video encoder with natural-language
summaries. Despite this activity, three fundamental gaps remain.

\paragraph{Gap 1: fragmented memory modalities.}
Each existing method targets a single memory type (spatial, semantic,
keyframe, feature). Combining them yields ad-hoc stacks with no
shared interface. A practitioner cannot cleanly ``turn on'' long-term
language memory in Mem-0 without rewriting its planner.

\paragraph{Gap 2: episode-bound memory.}
All published methods reset memory at episode end, precluding
``did I already put the knife in the top drawer yesterday?'' style
queries. Cross-session persistent memory---routine in LLM agents via
RAG---has not been integrated into VLA stacks.

\paragraph{Gap 3: incomparable benchmarks.}
MemoryVLA reports 96.5\% on LIBERO-5 but 1.5\% on the extended
MemoryBench used by ReMem-VLA. RMBench, MemoryBench, MIKASA-Robo, and
RoboCerebra each carve out a different subset of memory types with
different evaluation protocols. Cross-paper comparisons are
unreliable.

\subsection{Contributions.}
We introduce \method, a unified three-tier memory framework that
addresses all three gaps. Our design separates memory by time scale
and modality (\emph{hot} video, \emph{cold} language, \emph{persistent}
retrieval) and exposes them through a single planner prompt. We ship:

\begin{enumerate}
  \item A \textbf{semi-structured language memory schema}
        (Section~\ref{sec:method}) that compresses semantic events
        into human-readable text round-trippable with an LLM---the
        first fully interpretable long-term memory for VLAs.
  \item A \textbf{persistent cross-session store} with FAISS-ready
        retrieval and pluggable embedders, shippable as an inference
        hook for any existing VLA planner.
  \item A \textbf{concrete integration with Mem-0} \cite{rmbench}:
        only two code edits are required (Section~\ref{sec:method}),
        making adoption trivial for the broader community.
  \item A \textbf{synthetic retrieval benchmark}
        (Section~\ref{sec:experiments}) that verifies cross-session
        recall using only CPU, and a detailed proposal for
        \textbf{MNEMO-Bench}, a unified harness that consolidates
        RMBench, MemoryBench, MIKASA-Robo, and RoboCerebra into a
        single evaluation protocol with per-memory-type radar scores.
\end{enumerate}

Our release includes an MIT-licensed Python package
(\url{https://github.com/asimfish/mnemo-vla}) with 15 passing unit
tests on CPU, plus comprehensive documentation covering the
three-tier architecture, the Mem-0 integration, and the MNEMO-Bench
design. The full VLA training experiments are deferred to a follow-up
submission once the hot-memory tier is implemented; this paper
focuses on establishing the primitives, interfaces, and evaluation
methodology needed to make those experiments meaningful.
